\GAGASection{Appendix}{gaga:appendix}

All rings considered below are assumed to be \emph{commutative} and to have an
\emph{identity element}; all modules over these rings are assumed to be
\emph{unitary}.

\GAGASubsection{21. Flat modules}{gaga:no:21}

\GAGATarget{gaga:def:3}\textsc{Definition 3}. --- \emph{Let $B$ be an
$A$-module. We say that $B$ is $A$-flat (or flat) if, for every exact sequence
of $A$-modules}

\begin{equation*}
E\longrightarrow F\longrightarrow G,
\end{equation*}

\emph{the sequence}

\begin{equation*}
E\otimes_AB\longrightarrow F\otimes_AB\longrightarrow G\otimes_AB
\end{equation*}

\emph{is exact.}

By the definition of the Tor functors, the preceding condition is equivalent
to $\mathrm{Tor}_1^A(B,Q)=0$ for every $A$-module $Q$. Since Tor commutes with
inductive limits, it is enough to consider modules $Q$ of finite type, and
even, by the exact sequence for Tor, modules $Q$ generated by one element.
Thus $B$ is $A$-flat if and only if
$\mathrm{Tor}_1^A(B,A/\mathfrak{a})=0$ for every ideal $\mathfrak{a}$ of $A$;
equivalently, the canonical homomorphism
$\mathfrak{a}\otimes_AB\to B$ is injective.

\emph{Examples}. --- 1. If $A$ is a principal ideal domain, it follows from
the preceding discussion that `$B$ is $A$-flat' is equivalent to `$B$ is
torsion-free'.

2. If $S$ is a multiplicatively closed subset of a ring $A$, the ring of
fractions $A_S$ is $A$-flat, by \GAGACite{18}, no.~48, Lemma~1.

Let $A$ and $B$ be two rings, and let $\theta:A\to B$ be a homomorphism. This
homomorphism gives $B$ the structure of an $A$-module. If $E$ and $F$ are two
$A$-modules, then $E\otimes_AB$ and $F\otimes_AB$ have natural $B$-module
structures. Moreover, if $f:E\to F$ is a homomorphism, then $f\otimes1$ is a
$B$-homomorphism from $E\otimes_AB$ to $F\otimes_AB$. We thus obtain a
canonical $A$-linear map

\begin{equation*}
\mathrm{Hom}_A(E,F)\longrightarrow
\mathrm{Hom}_B(E\otimes_AB,F\otimes_AB),
\end{equation*}

which extends by linearity to a $B$-linear map

\begin{equation*}
\iota:\mathrm{Hom}_A(E,F)\otimes_AB\longrightarrow
\mathrm{Hom}_B(E\otimes_AB,F\otimes_AB).
\end{equation*}

\GAGATarget{gaga:prop:21}\textsc{Proposition 21}. --- \emph{The homomorphism
$\iota$ defined above is bijective when $A$ is a noetherian ring, $E$ is an
$A$-module of finite type, and $B$ is $A$-flat.}

For a fixed module $F$, put

\begin{equation*}
T(E)=\mathrm{Hom}_A(E,F)\otimes_AB
\qquad\text{and}\qquad
T'(E)=\mathrm{Hom}_B(E\otimes_AB,F\otimes_AB),
\end{equation*}

so that $\iota$ is a homomorphism from the functor $T(E)$ to the functor
$T'(E)$.

For $E=A$, we have $T(E)=T'(E)=F\otimes_AB$, and $\iota$ is bijective. The
same is true when $E$ is a free module of finite type.

Since $A$ is noetherian and $E$ is of finite type, there is an exact sequence

\begin{equation*}
L_1\longrightarrow L_0\longrightarrow E\longrightarrow0,
\end{equation*}

where $L_0$ and $L_1$ are free modules of finite type. Consider the
commutative diagram

\begin{center}
\begin{tikzcd}[column sep=large]
0 \arrow[r] & T(E) \arrow[r] \arrow[d,"\iota"'] &
T(L_0) \arrow[r] \arrow[d,"\iota_0"'] &
T(L_1) \arrow[d,"\iota_1"'] \\
0 \arrow[r] & T'(E) \arrow[r] & T'(L_0) \arrow[r] & T'(L_1).
\end{tikzcd}
\end{center}

The first row is exact because $B$ is $A$-flat; the second is exact by the
general properties of the functors $\otimes$ and $\mathrm{Hom}$. Since
$\iota_0$ and $\iota_1$ are bijective, it follows that $\iota$ is bijective,
as required.

\GAGASubsection{22. Flat pairs}{gaga:no:22}

\GAGATarget{gaga:def:4}\textsc{Definition 4}. --- \emph{Let $A$ be a ring and
let $B$ be a ring containing $A$. The pair $(A,B)$ is called flat if the
$A$-module $B/A$ is $A$-flat.}

We have:

\GAGATarget{gaga:prop:22}\textsc{Proposition 22}. --- \emph{A pair $(A,B)$ is
flat if and only if $B$ is $A$-flat and one of the following properties
holds:}

\GAGATarget{gaga:prop:22:item:a}\emph{a) (respectively a$'$)) For every
$A$-module (respectively every $A$-module of finite type) $E$, the
homomorphism $E\to E\otimes_AB$ is injective.}

\GAGATarget{gaga:prop:22:item:aa}\emph{a$''$) For every ideal $\mathfrak{a}$
of $A$, one has $\mathfrak{a}B\cap A=\mathfrak{a}$.}

If $E$ is any $A$-module, the exact sequence

\begin{equation*}
0\longrightarrow A\longrightarrow B\longrightarrow B/A\longrightarrow0
\end{equation*}

gives rise to the exact sequence

\begin{equation*}
\mathrm{Tor}_1^A(A,E)\longrightarrow
\mathrm{Tor}_1^A(B,E)\longrightarrow
\mathrm{Tor}_1^A(B/A,E)\longrightarrow
A\otimes_AE\longrightarrow B\otimes_AE.
\end{equation*}

Since $A\otimes_AE=E$ and $\mathrm{Tor}_1^A(A,E)=0$, we obtain the exact
sequence

\begin{equation*}
0\longrightarrow\mathrm{Tor}_1^A(B,E)\longrightarrow
\mathrm{Tor}_1^A(B/A,E)\longrightarrow E\longrightarrow E\otimes_AB.
\end{equation*}

Thus $\mathrm{Tor}_1^A(B/A,E)=0$ if and only if
$\mathrm{Tor}_1^A(B,E)=0$ and the homomorphism $E\to E\otimes_AB$ is
injective. The proposition follows immediately (observe that property
\hyperref[gaga:prop:22:item:aa]{a$''$)} says precisely that the homomorphism
$A/\mathfrak{a}\to A/\mathfrak{a}\otimes_AB$ is injective).

\GAGATarget{gaga:prop:23}\textsc{Proposition 23}. --- \emph{Let
$A\subset B\subset C$ be three rings. If the pairs $(A,C)$ and $(B,C)$ are
flat, then the pair $(A,B)$ is flat.}

We first show that $B$ is $A$-flat. In other words, given an exact sequence of
$A$-modules

\begin{equation*}
0\longrightarrow E\longrightarrow F,
\end{equation*}

we must show that
$0\to E\otimes_AB\to F\otimes_AB$ is exact. Let $N$ be the kernel of the
homomorphism $E\otimes_AB\to F\otimes_AB$. Since $C$ is $B$-flat, there is
an exact sequence

\begin{equation*}
0\longrightarrow N\otimes_BC\longrightarrow
(E\otimes_AB)\otimes_BC\longrightarrow(F\otimes_AB)\otimes_BC.
\end{equation*}

By associativity of the tensor product, $(E\otimes_AB)\otimes_BC$ is
identified with $E\otimes_AC$, and similarly $(F\otimes_AB)\otimes_BC$ is
identified with $F\otimes_AC$. Since $C$ is $A$-flat, the homomorphism
$E\otimes_AC\to F\otimes_AC$ is injective. It follows that $N\otimes_BC=0$;
applying Proposition~\hyperref[gaga:prop:22]{22} to the pair $(B,C)$, we see
that $N=0$. This proves that $B$ is $A$-flat.

On the other hand, if $E$ is any $A$-module, the composite homomorphism

\begin{equation*}
E\longrightarrow E\otimes_AB\longrightarrow E\otimes_AC
\end{equation*}

is injective because the pair $(A,C)$ is flat. Hence, \emph{a fortiori},
$E\to E\otimes_AB$ is injective. Thus the pair $(A,B)$ satisfies all the
hypotheses of Proposition~\hyperref[gaga:prop:22]{22}, as required.

\GAGATarget{gaga:prop:23:remark}\emph{Remark}. --- A similar argument shows
that if $(A,B)$ and $(B,C)$ are flat, then $(A,C)$ is flat. On the other hand,
it may happen that $(A,B)$ and $(A,C)$ are flat while $(B,C)$ is not.

\GAGASubsection{23. Modules over a local ring}{gaga:no:23}

Throughout this number, $A$ denotes a \emph{noetherian local ring}
\footnote{In fact, all the results proved in these last two numbers remain
valid without change for a \emph{Zariski ring} (cf. \GAGACite{15}, p.~157).} with
\emph{maximal ideal} $\mathfrak{m}$.

\GAGATarget{gaga:prop:24}\textsc{Proposition 24}. --- \emph{If an $A$-module
$E$ of finite type satisfies $E=\mathfrak{m}E$, then $E=0$.}

(See, for example, \GAGACite{15}, p.~138, or \GAGACite{4}, Exposé~I.)

Suppose $E\ne0$, and let $e_1,\ldots,e_n$ be a system of generators of $E$
having the smallest possible number of elements. Since
$e_n\in\mathfrak{m}E$, we have
$e_n=x_1e_1+\cdots+x_ne_n$, with $x_i\in\mathfrak{m}$, and hence

\begin{equation*}
(1-x_n)e_n=x_1e_1+\cdots+x_{n-1}e_{n-1}.
\end{equation*}

Since $1-x_n$ is invertible in $A$, this shows that
$e_1,\ldots,e_{n-1}$ generate $E$, contradicting the choice of $n$.

\GAGATarget{gaga:prop:24:cor}\textsc{Corollary}. --- \emph{Let $E$ be an
$A$-module of finite type. If a submodule $F$ of $E$ satisfies
$E=F+\mathfrak{m}E$, then $E=F$.}

Indeed, this relation means that $E/F=\mathfrak{m}(E/F)$.

We give every $A$-module $E$ the \emph{$\mathfrak{m}$-adic topology}, in which
the submodules $\mathfrak{m}^nE$ form a neighbourhood basis of $0$ (cf. \GAGACite{15},
p.~153).

\GAGATarget{gaga:prop:25}\textsc{Proposition 25}. --- \emph{Let $E$ be an
$A$-module of finite type. Then:}

\GAGATarget{gaga:prop:25:item:a}\emph{a) The topology induced on a submodule
$F$ of $E$ by the $\mathfrak{m}$-adic topology of $E$ coincides with the
$\mathfrak{m}$-adic topology of $F$.}

\GAGATarget{gaga:prop:25:item:b}\emph{b) Every submodule of $E$ is closed in
the $\mathfrak{m}$-adic topology of $E$ (and, in particular, $E$ is
separated).}

(See \GAGACite{15}, \emph{loc. cit.}, and \GAGACite{3}, Exposé~VIII bis.)

We briefly recall the proof. One begins by proving
\hyperref[gaga:prop:25:item:a]{a)}, either by using the theory of primary
decomposition (Krull; cf. \GAGACite{15}) or by establishing the existence of an
integer $r$ such that

\begin{equation*}
F\cap\mathfrak{m}^nE
=\mathfrak{m}^{n-r}(F\cap\mathfrak{m}^rE)
\quad\text{for}\quad n\geqslant r
\quad\text{(Artin--Rees; cf. \GAGACite{4}, Exposé~2).}
\end{equation*}

One then shows that $E$ is separated. Apply
\hyperref[gaga:prop:25:item:a]{a)} to the submodule $F$ which is the closure of
$0$ in $E$. Then $F=\mathfrak{m}F$, and hence $F=0$ by
Proposition~\hyperref[gaga:prop:24]{24}. Applying this result to quotient
modules of $E$ gives \hyperref[gaga:prop:25:item:b]{b)}.

Again let $E$ be an $A$-module of finite type, and let $\widehat{E}$ and
$\widehat{A}$ be the completions of $E$ and $A$ for the $\mathfrak{m}$-adic
topology. The bilinear map $A\times E\to E$ extends continuously to a map
$\widehat{A}\times\widehat{E}\to\widehat{E}$, making $\widehat{E}$ a
$\widehat{A}$-module. The canonical injection $E\to\widehat{E}$ therefore
extends by linearity to a homomorphism

\begin{equation*}
\varepsilon:E\otimes_A\widehat{A}\longrightarrow\widehat{E}.
\end{equation*}

\GAGATarget{gaga:prop:26}\textsc{Proposition 26}. --- \emph{For every
$A$-module $E$ of finite type, the homomorphism $\varepsilon$ defined above is
bijective.}

Let

\begin{equation*}
0\longrightarrow R\longrightarrow L\longrightarrow E\longrightarrow0
\end{equation*}

be an exact sequence of $A$-modules, with $L$ free of finite type. Since $A$
is noetherian, $R$ is of finite type. Proposition~\hyperref[gaga:prop:25]{25}
shows that the $\mathfrak{m}$-adic topology of $R$ is induced by that of $L$,
and it is clear that the topology of $E$ is the quotient topology of that of
$L$. Since these topologies are \emph{metrizable}, we obtain an exact sequence

\begin{equation*}
0\longrightarrow\widehat{R}\longrightarrow\widehat{L}
\longrightarrow\widehat{E}\longrightarrow0.
\end{equation*}

Consider the commutative diagram

\begin{center}
\begin{tikzcd}[column sep=large]
R\otimes_A\widehat{A} \arrow[r] \arrow[d,"\varepsilon''"'] &
L\otimes_A\widehat{A} \arrow[r] \arrow[d,"\varepsilon'"'] &
E\otimes_A\widehat{A} \arrow[r] \arrow[d,"\varepsilon"'] & 0 \\
\widehat{R} \arrow[r] & \widehat{L} \arrow[r] & \widehat{E} \arrow[r] & 0.
\end{tikzcd}
\end{center}

The two rows are exact, and $\varepsilon'$ is plainly bijective. It follows
that $\varepsilon$ is surjective (that is,
$\widehat{E}=\widehat{A}\mathbin{.}E$; cf. \GAGACite{15}, p.~153, Lemma~1). Since this
has been proved for every $A$-module of finite type, it applies in particular
to $R$, showing that $\varepsilon''$ is surjective. The five lemma now shows
that $\varepsilon$ is bijective, as required.

\GAGASubsection{24. Flatness properties of local rings}{gaga:no:24}

All local rings considered below are assumed to be noetherian.

\GAGATarget{gaga:prop:27}\textsc{Proposition 27}. --- \emph{Let $A$ be a
local ring and let $\widehat{A}$ be its completion. The pair
$(A,\widehat{A})$ is flat.}

First, $\widehat{A}$ is $A$-flat. Indeed, it is enough to show that, if
$E\to F$ is injective, then so is
$E\otimes_A\widehat{A}\to F\otimes_A\widehat{A}$, and we may even assume
that $E$ and $F$ are of finite type. In that case,
Proposition~\hyperref[gaga:prop:26]{26} identifies
$E\otimes_A\widehat{A}$ with $\widehat{E}$ and
$F\otimes_A\widehat{A}$ with $\widehat{F}$; the assertion follows from the
evident fact that $\widehat{E}$ embeds in $\widehat{F}$.

Similarly, the injectivity of $E\to\widehat{E}$ for $E$ of finite type shows
that the pair $(A,\widehat{A})$ satisfies property
\hyperref[gaga:prop:22:item:a]{a$'$)} of
Proposition~\hyperref[gaga:prop:22]{22}; hence it is a flat pair.

Now let $A$ and $B$ be two local rings, and let $\theta:A\to B$ be a
homomorphism. Suppose that $\theta$ maps the maximal ideal of $A$ into the
maximal ideal of $B$. Then $\theta$ is continuous and extends continuously to
a homomorphism $\widehat{\theta}:\widehat{A}\to\widehat{B}$.

\GAGATarget{gaga:prop:28}\textsc{Proposition 28}. --- \emph{Suppose that
$\widehat{\theta}:\widehat{A}\to\widehat{B}$ is bijective, and identify $A$
with a subring of $B$ by means of $\theta$. Then the pair $(A,B)$ is flat.}

We have $A\subset B\subset\widehat{B}=\widehat{A}$, and the pairs
$(A,\widehat{A})$ and $(B,\widehat{B})$ are flat by the preceding proposition.
Proposition~\hyperref[gaga:prop:23]{23} shows that $(A,B)$ is a flat pair.

\GAGATarget{gaga:prop:29}\textsc{Proposition 29}. --- \emph{Let $A$ and $B$
be two local rings, let $\mathfrak{a}$ be an ideal of $A$, and let
$\theta:A\to B$ be a homomorphism. If $\theta$ satisfies the hypothesis of
Proposition~\hyperref[gaga:prop:28]{28}, then so does the homomorphism from
$A/\mathfrak{a}$ to $B/\theta(\mathfrak{a})B$ defined by $\theta$. Thus the
pair $(A/\mathfrak{a},B/\theta(\mathfrak{a})B)$ is flat.}

By Proposition~\hyperref[gaga:prop:26]{26}, the completion of
$A/\mathfrak{a}$ is $\widehat{A}/\mathfrak{a}\widehat{A}$; similarly, the
completion of $B/\theta(\mathfrak{a})B$ is
$\widehat{B}/\theta(\mathfrak{a})\widehat{B}$. The result follows.

\GAGATarget{gaga:prop:30}\textsc{Proposition 30}. --- \emph{Let $A$ and $A'$
be two local rings, let $\theta:A\to A'$ be a homomorphism satisfying the
hypothesis of Proposition~\hyperref[gaga:prop:28]{28}, and let $E$ be an
$A$-module of finite type. If the $A'$-module
$E'=E\otimes_AA'$ is isomorphic to ${A'}^n$, then $E$ is isomorphic to
$A^n$.}

We identify $A$ with a subring of $A'$ by means of $\theta$. If
$\mathfrak{m}$ and $\mathfrak{m}'$ denote the maximal ideals of $A$ and $A'$,
then $\mathfrak{m}\subset\mathfrak{m}'$. Moreover, since $\mathfrak{m}'$ is a
neighbourhood of $0$ in $A'$ and $A$ is dense in $A'$, we have
$A'=\mathfrak{m}'+A$. It follows that

\begin{equation*}
A/\mathfrak{m}=A'/\mathfrak{m}',
\qquad
E/\mathfrak{m}E=E'/\mathfrak{m}'E'.
\end{equation*}

Since the $A'$-module $E'$ is free of rank $n$, the
$A'/\mathfrak{m}'$-module $E'/\mathfrak{m}'E'$ is free of rank $n$. We can
therefore choose $n$ elements $e_1,\ldots,e_n$ of $E$ whose images in
$E/\mathfrak{m}E$ form a basis of that vector space over
$A/\mathfrak{m}$. The elements $e_i$ define a homomorphism
$f:A^n\to E$ which is surjective by the corollary to
Proposition~\hyperref[gaga:prop:24]{24}. We shall prove that $f$ is injective,
which will prove the proposition.

Let $N$ be the kernel of $f$. Since the pair $(A,A')$ is flat
(Proposition~\hyperref[gaga:prop:28]{28}), the exact sequence

\begin{equation*}
0\longrightarrow N\longrightarrow A^n\xrightarrow{f}E\longrightarrow0
\end{equation*}

gives rise to the exact sequence

\begin{equation*}
0\longrightarrow N'\longrightarrow {A'}^n
\xrightarrow{f'}E'\longrightarrow0.
\end{equation*}

Since $E'$ is free, $N'$ is a direct summand of ${A'}^n$, and there is an
exact sequence

\begin{equation*}
0\longrightarrow N'/\mathfrak{m}'N'
\longrightarrow {A'}^n/\mathfrak{m}'{A'}^n
\longrightarrow E'/\mathfrak{m}'E'\longrightarrow0.
\end{equation*}

By construction, $f'$ induces a bijection from
${A'}^n/\mathfrak{m}'{A'}^n$ onto $E'/\mathfrak{m}'E'$. Hence
$N'/\mathfrak{m}'N'=0$, so $N'=0$ by
Proposition~\hyperref[gaga:prop:24]{24}; and then $N=0$ because the pair
$(A,A')$ is flat, as required.

\GAGASection{Bibliography}{gaga:bibliography}

\begin{enumerate}[label={[\arabic*]},leftmargin=3em]
\item\GAGATarget{gaga:bib:1} H.~\textsc{Cartan}. Idéaux et modules de fonctions analytiques de
variables complexes. \emph{Bull. Soc. Math. France}, \textbf{78}, 1950,
pp.~29--64.

\item\GAGATarget{gaga:bib:2} H.~\textsc{Cartan}. Séminaire E. N. S., 1951--1952.

\item\GAGATarget{gaga:bib:3} H.~\textsc{Cartan}. Séminaire E. N. S., 1953--1954.

\item\GAGATarget{gaga:bib:4} H.~\textsc{Cartan} and C.~\textsc{Chevalley}. Séminaire E. N. S.,
1955--1956.

\item\GAGATarget{gaga:bib:5} H.~\textsc{Cartan} and J.-P.~\textsc{Serre}. Un théorème de finitude
concernant les variétés analytiques compactes. \emph{C. R.}, \textbf{237},
1953, pp.~128--130.

\item\GAGATarget{gaga:bib:6} W-L.~\textsc{Chow}. On compact complex analytic varieties.
\emph{Amer. J. of Maths.}, \textbf{71}, 1949, pp.~893--914.

\item\GAGATarget{gaga:bib:7} W-L.~\textsc{Chow}. On the projective embedding of homogeneous
varieties. \emph{Lefschetz's volume}, Princeton, 1956.

\item\GAGATarget{gaga:bib:8} P.~\textsc{Dolbeault}. Sur la cohomologie des variétés analytiques
complexes. \emph{C. R.}, \textbf{236}, 1953, pp.~175--177.

\item\GAGATarget{gaga:bib:9} J.~\textsc{Frenkel}. Cohomologie à valeurs dans un faisceau non
abélien. \emph{C. R.}, \textbf{240}, 1955, pp.~2368--2370.

\item\GAGATarget{gaga:bib:10} A.~\textsc{Grothendieck}. A general theory of fibre spaces with
structure sheaf. \emph{Kansas Univ.}, 1955.

\item\GAGATarget{gaga:bib:11} K.~\textsc{Kodaira}. On Kähler varieties of restricted type (an
intrinsic characterization of algebraic varieties). \emph{Ann. of Maths.},
\textbf{60}, 1954, pp.~28--48.

\item\GAGATarget{gaga:bib:12} K.~\textsc{Kodaira} and D. C.~\textsc{Spencer}. Divisor class groups
on algebraic varieties. \emph{Proc. Nat. Acad. Sci. U. S. A.}, \textbf{39},
1953, pp.~872--877.

\item\GAGATarget{gaga:bib:13} M.~\textsc{Rosenlicht}. Some basic theorems on algebraic groups.
\emph{Amer. J. of Maths.}, \textbf{78}, 1956, pp.~401--443.

\item\GAGATarget{gaga:bib:14} W.~\textsc{Rückert}. Zum Eliminationsproblem der Potenzreihenideale.
\emph{Math. Ann.}, \textbf{107}, 1933, pp.~259--281.

\item\GAGATarget{gaga:bib:15} P.~\textsc{Samuel}. Commutative Algebra (Notes by D. Herzig).
\emph{Cornell Univ.}, 1953.

\item\GAGATarget{gaga:bib:16} P.~\textsc{Samuel}. Algèbre locale. \emph{Mém. Sci. Math.}, 123,
Paris, 1953.

\item\GAGATarget{gaga:bib:17} P.~\textsc{Samuel}. Méthodes d'algèbre abstraite en géométrie
algébrique. \emph{Ergebn. der Math.}, Springer, 1955.

\item\GAGATarget{gaga:bib:18} J.-P.~\textsc{Serre}. Faisceaux algébriques cohérents.
\emph{Ann. of Maths.}, \textbf{61}, 1955, pp.~197--278.

\item\GAGATarget{gaga:bib:19} J.-P.~\textsc{Serre}. Sur la cohomologie des variétés algébriques.
\emph{J. de Maths. Pures et Appl.}, \textbf{35}, 1956.

\item\GAGATarget{gaga:bib:20} A.~\textsc{Weil}. Fibre-spaces in algebraic geometry (Notes by
A. Wallace). \emph{Chicago Univ.}, 1952.
\end{enumerate}

\begin{center}\rule{4cm}{0.4pt}\end{center}
